\chapter{General Architecture}
\label{ch:architecture}

\section{Two phases: indexing and mapping}

shmap's execution splits into two clearly separated phases, run once per invocation:

\begin{enumerate}
  \item \textbf{Indexing phase} (Chapter~\ref{ch:indexing}): read the reference FASTA once,
    FracMinHash-sketch every segment, and build a hash map from $k$-mer hash to the list of
    reference positions (\emph{hits}) sharing that hash. This is done once regardless of how many
    reads are mapped afterward.
  \item \textbf{Mapping phase} (Chapters~\ref{ch:seeding}--\ref{ch:scoring},
    reassembled end-to-end in Chapter~\ref{ch:end-to-end}): stream the read FASTA/FASTQ one record
    at a time; for each read, sketch it, seed it against the index, bucket the hits, prune buckets
    with the seed heuristic, refine survivors, compute MAPQ from the best two survivors, and emit
    one PAF line.
\end{enumerate}

\begin{figure}[htbp]
\centering
\begin{tikzpicture}[
    node distance=7mm and 12mm,
    every node/.style={font=\small},
    box/.style={draw, rounded corners, fill=blue!6, minimum height=9mm, text width=32mm, align=center, thick},
    dbox/.style={draw, rounded corners, fill=orange!10, minimum height=9mm, text width=32mm, align=center, thick},
    io/.style={draw, dashed, fill=gray!5, minimum height=8mm, text width=28mm, align=center},
    arr/.style={-{Stealth[length=2.2mm]}, thick}
  ]
  % Indexing phase (left column)
  \node[io] (reffa) {Reference FASTA};
  \node[box, below=of reffa] (sketchref) {FracMinHash sketch\\ each segment (\S\ref{ch:sketching})};
  \node[box, below=of sketchref] (index) {Build \code{h2single}/\code{h2multi}\\ hash$\to$hits index (\S\ref{ch:indexing})};
  \node[io, below=of index] (idx) {SketchIndex};

  \draw[arr] (reffa) -- (sketchref);
  \draw[arr] (sketchref) -- (index);
  \draw[arr] (index) -- (idx);

  % Mapping phase (right column), per read
  \node[io, right=28mm of reffa] (readfa) {Read FASTA/FASTQ\\ (one record)};
  \node[box, below=of readfa] (sketchread) {FracMinHash sketch\\ the read (\S\ref{ch:sketching})};
  \node[box, below=of sketchread] (seed) {Group into seeds,\\ rarest-first (\S\ref{ch:seeding})};
  \node[box, below=of seed] (bucket) {Match seeds into\\ overlapping buckets (\S\ref{ch:bucketing})};
  \node[dbox, below=of bucket] (prune) {Seed-heuristic pruning\\ (\S\ref{ch:pruning})};
  \node[box, below=of prune] (score) {Refine survivors:\\ Containment/Jaccard/\\bucket\_SH/bucket\_LCS (\S\ref{ch:scoring})};
  \node[box, below=of score] (mapq) {Best + 2nd-best $\to$ MAPQ\\ (\S\ref{ch:mapq})};
  \node[io, below=of mapq] (paf) {PAF line (\S\ref{ch:output})};

  \draw[arr] (readfa) -- (sketchread);
  \draw[arr] (sketchread) -- (seed);
  \draw[arr] (seed) -- (bucket);
  \draw[arr] (bucket) -- (prune);
  \draw[arr] (prune) -- (score);
  \draw[arr] (score) -- (mapq);
  \draw[arr] (mapq) -- (paf);

  \draw[arr] (idx.east) to[out=0,in=180] (bucket.west);
  \node at ($(idx)!0.5!(bucket)+(0,0)$) {};

  \node[draw, dotted, fit=(reffa)(idx), inner sep=3mm, label={[font=\bfseries\small]above:Indexing phase (once)}] (fitL) {};
  \node[draw, dotted, fit=(readfa)(paf), inner sep=3mm, label={[font=\bfseries\small]above:Mapping phase (per read)}] (fitR) {};
\end{tikzpicture}
\caption{shmap's two-phase pipeline. The indexing phase runs once; the mapping phase repeats for
every read in the query file, consulting the already-built index.}
\label{fig:pipeline}
\end{figure}

\section{Configuration axes: the three compile-time flags}
\label{sec:config-axes}

The reference C++ implementation parameterizes its core mapper class as a template,
\code{SHMapper<no\_bucket\_pruning, one\_sweep, abs\_pos>}, over three independent booleans. In
the original source these are \emph{compile-time} template parameters (so that the compiler can
specialize/inline branch-free code for whichever combination is selected at build time), but they
are semantically just three independent feature toggles, exposed at the CLI as \code{-P}, \code{-B},
\code{-F} respectively:

\begin{center}
\begin{tabular}{@{}lp{9.5cm}@{}}
\toprule
\textbf{Flag / parameter} & \textbf{Effect} \\
\midrule
\code{no\_bucket\_pruning} (\code{-P}) &
  When true, \emph{disables} the seed-heuristic pruning of \S\ref{sec:seed-heuristic-theory}
  entirely: every candidate bucket is fully scored. Correct but much slower; useful as a
  ground-truth oracle to verify pruning never discards a bucket it shouldn't. \\
\code{one\_sweep} (\code{-B}) &
  Intended (per its own \texttt{-B} help text, ``disregards seed heuristic and runs one sweep on
  all matches'') to skip the incremental per-bucket seed budget entirely and consider every seed
  at once. \emph{Chapter~\ref{ch:defects} documents that in the reference source this flag is
  fully dead: parsed, templated, and printed, but never once branched on inside the mapper.} \\
\code{abs\_pos} (\code{-F}) &
  Selects the coordinate space used for bucket indexing and window-size comparisons: \emph{sketch
  index} (position within the array of sketch $k$-mers, i.e. ``the $i$-th kept $k$-mer'') when
  false (default), or \emph{absolute nucleotide position} when true. This changes bucket boundary
  arithmetic throughout Chapters~\ref{ch:bucketing}--\ref{ch:scoring} (every place that reads
  \emph{``in sketch-index units, or nucleotide units if \code{abs\_pos}''}). \\
\bottomrule
\end{tabular}
\end{center}

\begin{notebox}
A from-scratch implementation does not need compile-time specialization to get shmap's
performance characteristics: these three booleans can equally well be ordinary runtime
parameters threaded through the mapping functions (this is exactly what the accompanying Rust
re-implementation does, using const generics only to keep the option of zero-cost specialization
available, not because runtime dispatch is fundamentally required). What \emph{is} essential to
preserve is that each one is read consistently by \emph{every} function that needs it --- the
original source's own bug catalogue (Chapter~\ref{ch:defects}) includes at least one flag that
looks load-bearing but silently isn't.
\end{notebox}

\section{The metric axis}
\label{sec:metric-axis}

Orthogonal to the three structural flags above, a fourth, purely-scoring-level axis selects
\emph{how} a surviving bucket's similarity score is computed (CLI flag \code{-m}, four values):
\code{Containment}, \code{Jaccard}, \code{bucket\_SH}, \code{bucket\_LCS}. These map directly onto
the theory of \S\ref{sec:fracminhash-theory} (containment/Jaccard) and \S\ref{sec:lis-theory}
(LCS); \code{bucket\_SH} is a degenerate/cheap option that reuses the seed-heuristic bound itself
as the final score (no refinement sweep at all). Full formulas are in
\S\ref{sec:metric-variants}.

\section{Data-flow summary: what crosses each module boundary}

\begin{center}
\begin{tabular}{@{}p{3.1cm}p{4.6cm}p{6cm}@{}}
\toprule
\textbf{Stage} & \textbf{Input} & \textbf{Output} \\
\midrule
Sketch reference & raw nucleotide bytes per segment & \code{Vec<Kmer>} per segment (position,
  hash, strand) \\
Build index & all segments' \code{Vec<Kmer>} & \code{hash} $\to$ \code{Hit} (unique) or
  \code{hash} $\to$ \code{Vec<Hit>} (repeated), each \code{Hit} carrying segment id + position \\
Sketch read & raw nucleotide bytes of the read & \code{Vec<Kmer>} \\
Seed & \code{Vec<Kmer>} of the read & \code{Vec<Seed>}: one entry per distinct hash, with
  occurrence count in the read and reference hit count, sorted rarest-reference-hit-first \\
Bucket / match seeds & \code{Vec<Seed>} (budget-limited prefix) + index & per-\code{(segment,
  bucket)} running \code{BucketContent} (seed count, match count, codirection, position range) \\
Prune & sorted buckets (most-matches-first) + remaining unconsumed seeds & surviving buckets only
  (most are discarded using $O(1)$ arithmetic per extension step) \\
Refine & one surviving bucket + the read's seed table & one \code{Mapping} (score, matched
  interval, strand, intersection count) \\
Best/second-best & every refined \code{Mapping} & the two highest, non-overlapping,
  \code{Mapping}s \\
MAPQ & best two scores + same-strand-seed counts & integer MAPQ in $\{0,5,60\}$ (shmap's concrete
  range; see Chapter~\ref{ch:mapq}) \\
Output & final \code{Mapping} (+ globally accumulated counters) & one PAF text line with custom
  tags \\
\bottomrule
\end{tabular}
\end{center}

\section{Complexity at a glance}

Let $n = |T|$ (total reference length), $m$ = a read's sketch size ($\approx r\cdot|P|$), and $S
\le m$ the seed budget actually consumed for that read (\S\ref{sec:seed-budget}). Indexing is
$O(n)$ time (one rolling-hash pass) and $O(rn)$ space (the kept fraction of $k$-mers). Per read:
seeding/sorting is $O(m\log m)$; seed matching against the index is $O(S \cdot \bar d)$ where
$\bar d$ is the average number of reference hits per consumed seed (bounded above by \code{-M},
the max-matches-per-kmer cap, \S\ref{sec:max-matches}); bucket pruning is, in the typical case
where most buckets are eliminated in $O(1)$ extensions, sublinear in the number of buckets touched;
and refinement of each surviving bucket costs $O(|\text{bucket}|)$ (a single linear sweep). Full
derivations are in Chapter~\ref{ch:complexity}.
